\documentclass[a4paper,11pt]{article}
        \usepackage[top=15mm,bottom=18mm,left=16mm,right=16mm]{geometry}
        \usepackage{fontspec}
        \usepackage{xeCJK}
        \setmainfont{Times New Roman}
        \setCJKmainfont{Yu Gothic}
        \setmonofont{Consolas}
        \setCJKmonofont{Yu Gothic}
        \usepackage{booktabs}
        \usepackage{longtable}
        \usepackage{tabularx}
        \usepackage{array}
        \usepackage{enumitem}
        \usepackage{hyperref}
        \usepackage{listings}
        \usepackage{xcolor}
        \usepackage{amsmath}
        \usepackage{amssymb}
        \hypersetup{
          colorlinks=true,
          linkcolor=blue,
          urlcolor=blue,
          pdftitle={profile YAML 読込と検証},
          pdfauthor={Codex}
        }
        \lstset{
          basicstyle=\ttfamily\small,
          breaklines=true,
          columns=fullflexible,
          keepspaces=true,
          frame=single,
          framerule=0.2pt,
          rulecolor=\color{black},
          showstringspaces=false
        }
        \setlength{\parskip}{0.42em}
        \setlength{\parindent}{1em}
        \setlength{\emergencystretch}{3em}
        \renewcommand{\arraystretch}{1.15}
        \title{08. profile YAML 読込と検証}
        \author{solar\_ws0129-main}
        \date{2026-07-29}
        \begin{document}
        \maketitle
        \tableofcontents
        \clearpage

        \section{対象}
        \begin{tabularx}{\linewidth}{>{\raggedright\arraybackslash}p{0.18\linewidth} X}
        \toprule
        項目 & 内容 \\
        \midrule
        ファイル & \path|mpc_solarcar/solar_profile.py| \\
        ソースSHA-256 & \path|342c5f001c7ae09a4ccfafe19e30706d124ac8b044852c365ca03ee664af00e7| \\
        種別 & Python \\
        区分 & config \\
        役割 & profile.yaml をロードし、セクション取得、相対パス解決、CSV 最低品質検査を行う。 \\
        起動文脈 & launch、offline simulation、identification の共通設定入口。 \\
        \bottomrule
        \end{tabularx}

        \section{このファイルを読む理由}
        profile.yaml をロードし、セクション取得、相対パス解決、CSV 最低品質検査を行う。

        \begin{itemize}[leftmargin=1.6em]
        \item load\_profile が YAML 全体を返す。
\item get\_path が profile 基準で実ファイルパスへ変換する。
\item require\_csv\_data\_rows が空テンプレと実データを区別する。
        \end{itemize}

        \section{呼び出し関係}
        \subsection{どこから呼ばれるか}
        \begin{itemize}[leftmargin=1.6em]
        \item \path|live_launch.py|
\item \path|solarcar_sim.launch.py|
\item \path|solar_state_node.py|
\item 多数の scripts
        \end{itemize}

        \subsection{次に読むと理解が進むファイル}
        \begin{itemize}[leftmargin=1.6em]
        \item \path|mpc_solarcar/path_utils.py|
        \end{itemize}

        \subsection{同一リポジトリ内の主要依存}
        \begin{itemize}[leftmargin=1.6em]
        \item 同一リポジトリ内の明示 import は少ない。
        \end{itemize}

        \section{ソース冒頭の把握}
        モジュール docstring または先頭意図の要約:

        モジュール docstring は明示されていない。

        \subsection{代表コード断片}
        \begin{lstlisting}[language=Python]
ROOT = Path(__file__).resolve().parents[1]


def require_csv_data_rows(
    path: str | Path,
    *,
    label: str,
    required_columns: tuple[str, ...] = (),
) -> Path:
    resolved = Path(path).expanduser().resolve()
    if not resolved.is_file():
        raise FileNotFoundError(f"{label} CSV was not found: {resolved}")
    with resolved.open("r", encoding="utf-8-sig", newline="") as stream:
        reader = csv.DictReader(stream)
        header = tuple(reader.fieldnames or ())
        missing = [column for column in required_columns if column not in header]
        if missing:
            raise ValueError(f"{label} CSV is missing columns {missing}: {resolved}")
        if next(reader, None) is None:
            raise ValueError(
                f"{label} CSV has a header but no data rows: {resolved}. "
                "Fill the template or select an identified vehicle profile before launch."
\end{lstlisting}


        \section{主要構造の読み取り}
        主要関数は require\_csv\_data\_rows, resolve\_relative\_path, load\_profile, merged\_dict, get\_section, get\_path。

        \subsection{主要クラス・関数}
            \begin{longtable}{p{0.07\linewidth} p{0.16\linewidth} p{0.25\linewidth} p{0.40\linewidth}}
            \toprule
            行 & 種別 & 名前 & 補足 \\
            \midrule
            14 & function & \path|require_csv_data_rows| & args: path \\
37 & function & \path|_resolve_profile_path| & args: path\_like \\
51 & function & \path|resolve_relative_path| & args: base\_dir, path\_like \\
68 & function & \path|load_profile| & args: path\_like \\
77 & function & \path|merged_dict| &  \\
78 & function & \path|_merge| & args: dst, src \\
93 & function & \path|get_section| & args: cfg, key, default \\
108 & function & \path|get_path| & args: cfg, profile\_path, key, default \\
            \bottomrule
            \end{longtable}

        \subsection{ファイルを上から読んだときの定義順}
            \begin{longtable}{p{0.07\linewidth} p{0.84\linewidth}}
            \toprule
            行 & 内容 \\
            \midrule
            11 & ROOT に Path(\_\_file\_\_).resolve().parents[1] の結果を代入する。 \\
14 & 関数 require\_csv\_data\_rows を定義する。 \\
37 & 関数 \_resolve\_profile\_path を定義する。 \\
51 & 関数 resolve\_relative\_path を定義する。 \\
68 & 関数 load\_profile を定義する。 \\
77 & 関数 merged\_dict を定義する。 \\
93 & 関数 get\_section を定義する。 \\
108 & 関数 get\_path を定義する。 \\
            \bottomrule
            \end{longtable}

        \subsection{import 群の詳細}
            \begin{longtable}{p{0.07\linewidth} p{0.33\linewidth} p{0.50\linewidth}}
            \toprule
            行 & import 文 & このファイルで読む理由 \\
            \midrule
            1 & \path|from __future__ import annotations| & 型ヒントの遅延評価を有効にし、自己参照型や forward reference を安全に扱うため。 このファイル内での使用位置は少ないか、間接利用である。 \\
3 & \path|import copy| & 設定辞書や payload を安全に複製するため。 このファイル内での主な使用位置は L81, L83, L99, L102, L104, L105。 \\
4 & \path|import csv| & CSV の逐次読込・逐次書込を行うため。 このファイル内での主な使用位置は L24。 \\
5 & \path|from pathlib import Path| & ファイルやディレクトリを安全に扱うため。 このファイル内での主な使用位置は L11, L15, L19, L20, L37, L38, L42, L51, ...。 \\
6 & \path|from typing import Any| & typing から Any を読み込み、このファイルの処理を組み立てるため。 このファイル内での主な使用位置は L86, L93, L94。 \\
8 & \path|import yaml| & profile、stop、schedule、summary YAML を読み書きするため。 このファイル内での主な使用位置は L71。 \\
            \bottomrule
            \end{longtable}

        \section{このファイルを読む前に必要な基礎知識}
        次の章は、構文やROS用語を既知と仮定しないための説明である。

        \subsection{プログラム、プロセス、メモリ、オブジェクトを区別する}
ソースファイルはディスク上の文字列であり、それ自体は走っていない。Python実行可能プログラムがソースを読み、OSがその実行を一つのプロセスとして管理する。プロセスには仮想メモリ、開いているファイル、スレッド、終了コードなどが対応する。


\begin{lstlisting}
ソースファイル -> Pythonインタプリタが読み込む -> OS上のプロセス -> Pythonオブジェクトをメモリに生成 -> 関数やcallbackを実行
\end{lstlisting}


「メモリ上に生成する」とは、実行中プロセスが使う記憶領域に、その値の型、属性、参照関係を表す実体を用意することである。変数は実体そのものというより、そのオブジェクトを指す名前として理解するとPythonの代入が読みやすい。


\begin{lstlisting}[language=python]
node = MPCNode()
alias = node
# nodeとaliasは同じオブジェクトを参照する。
\end{lstlisting}


一つのプロセスには複数スレッドを持たせられる。スレッドは同じプロセスのメモリを共有するため、受信callbackと最適化callbackが同じself.zを書き換える場合は実行順と排他を検討する必要がある。

\paragraph{根拠資料}
\begin{itemize}[leftmargin=1.6em]
\item \href{https://docs.python.org/3/tutorial/classes.html}{Python公式チュートリアル: Classes}
\end{itemize}

\subsection{名前、代入、型、参照}
Pythonの代入文は、右辺を先に評価し、その結果のオブジェクトへ左辺の名前を結び付ける。`x = f()`では、まずfを呼び、その戻り値が得られてからxが更新される。


`float(...)`、`int(...)`、`bool(...)`は型名を呼び出して値を変換する。外部CSV、YAML、ROS parameterから来た値は型が期待どおりとは限らないため、このリポジトリでは明示変換が多い。


`None`は値が無いことを示す単一のオブジェクト、`math.nan`は浮動小数点値ではあるが有効な数値ではないことを示す。両者は用途が異なるため、`is None`と`math.isfinite`を使い分ける。

\paragraph{根拠資料}
\begin{itemize}[leftmargin=1.6em]
\item \href{https://docs.python.org/3/tutorial/classes.html}{Python公式チュートリアル: Classes}
\end{itemize}

\subsection{関数、引数、戻り値、スコープ}
`def`は関数本体をその場で実行する文ではない。関数オブジェクトを作り、その名前を現在の名前空間へ登録する。`f`は関数そのもの、`f()`はその関数を呼んだ結果である。


仮引数は関数定義側の受け取り名、実引数は呼出側から渡す値である。位置引数、キーワード引数、既定値付き引数、`*args`、`**kwargs`は渡し方が異なる。


\begin{lstlisting}[language=python]
def clip_speed(value, minimum=0.0, maximum=110.0):
    return min(maximum, max(minimum, value))

v = clip_speed(120.0, maximum=90.0)
\end{lstlisting}


関数内で作った通常の名前はローカルスコープに属する。メソッドが`self.z`を更新すればオブジェクトの状態が残るが、単に`z`を更新しただけならその呼出しのローカル値である。

\paragraph{根拠資料}
\begin{itemize}[leftmargin=1.6em]
\item \href{https://docs.python.org/3/tutorial/classes.html}{Python公式チュートリアル: Classes}
\end{itemize}

\subsection{module、package、import、console\_scripts}
Pythonファイルはmoduleとして読み込める。複数moduleをディレクトリにまとめたものがpackageである。`from .model import SolarCarModel`の先頭の点は、現在と同じpackage内のmodel moduleを指す相対importである。


import時にはファイルのトップレベル文が上から一度実行される。`def`や`class`は関数・クラスを登録するが、その本体の通常処理は呼び出すまで走らない。


setuptoolsの`console\_scripts`は、端末で使う実行可能名と`package.module:function`を対応付けるインストール時メタデータである。生成される実行用ラッパーはmoduleをimportし、指定関数を呼び、戻り値を終了コードとして扱う小さな入口である。


\begin{lstlisting}
ROS launchのexecutable名 -> install済みconsole script -> mpc_solarcar.mpc_nodeをimport -> main()を呼ぶ
\end{lstlisting}

\paragraph{根拠資料}
\begin{itemize}[leftmargin=1.6em]
\item \href{https://setuptools.pypa.io/en/latest/userguide/entry_point.html}{setuptools公式: Entry Points / Console Scripts}
\item \href{https://docs.python.org/3/tutorial/classes.html}{Python公式チュートリアル: Classes}
\end{itemize}

\subsection{例外、try/finally、with、資源解放}
例外は通常の戻り値とは別経路で異常を呼出元へ伝える。`try/except`は想定した異常を処理し、`finally`は成功・失敗にかかわらず後始末を行う。


`with open(...) as f:`はcontext managerを使い、ブロックを出るとファイルを閉じる。CSVやログの破損を避けるため、開いた資源の所有者と閉じる場所を明確にする。


`except Exception: pass`はノードを止めない利点がある一方、入力異常を隠して原因追跡を難しくする。安全に関係する値では、少なくとも頻度制限付き警告、異常カウンタ、fallback状態のpublishを検討する。



\subsection{制御、状態、入力、モデル予測制御MPC}
制御対象の内部を表す状態をx、操作入力をu、外乱・予報をwとすると、離散モデルは`x[k+1] = f(x[k], u[k], w[k])`と書ける。ソーラーカーではSoC、電池温度、距離、速度などが状態候補、速度目標や駆動トルクが入力候補になる。


\[
\min_{u_0,\ldots,u_{N-1}} \sum_{k=0}^{N-1}\ell(x_k,u_k,w_k)+V_f(x_N)
\]

MPCは現在状態からNステップ先まで予測し、目的関数と制約を満たす入力系列を求める。ただし実際に適用するのは通常先頭入力だけで、次回は新しい実測状態から再び解く。これがreceding horizonである。


予測モデル、目的関数、制約、ホライズン、solver、初期値のどれかが変わると答えも変わる。「MPCを使う」だけでは仕様は決まらず、これらを単位付きで追う必要がある。

\paragraph{根拠資料}
\begin{itemize}[leftmargin=1.6em]
\item \href{https://docs.scipy.org/doc/scipy/reference/generated/scipy.optimize.minimize.html}{SciPy公式: scipy.optimize.minimize}
\end{itemize}



        \section{関数・クラスを上から順に解説}
        \subsection{L14 関数 require\_csv\_data\_rows}
            \begin{itemize}[leftmargin=1.6em]
              \item 行範囲: L14--L34
              \item 所属: module直下
              \item 定義: \path!require_csv_data_rows(path: str | Path, *, label: str, required_columns: tuple[str, ...] = ()) -> Path!
              \item docstring: 明示されていない。
              \item このブロックが直接呼ぶ主な関数・メソッド: DictReader, FileNotFoundError, Path, ValueError, expanduser, is\_file, next, open, resolve, tuple
              \item 戻り値の要点: resolved
              \item 主なローカル代入名: column, header, missing, reader, resolved, stream
              \item 読み取る主なインスタンス属性: 特になし。
              \item 更新する主なインスタンス属性: 特になし。
              \item 明示的に送出する例外: FileNotFoundError(f'\{label\} CSV was not found: \{resolved\}'), ValueError(f'\{label\} CSV has a header but no data rows: \{resolved\}. Fill the template or select an identified vehicle profile before launch.'), ValueError(f'\{label\} CSV is missing columns \{missing\}: \{resolved\}')
              \item 制御構造の規模: 条件分岐 3、ループ 0、try 0
            \end{itemize}
            \paragraph{この定義を読むためのPython構文}
            \begin{itemize}[leftmargin=1.6em]
            \item `->`以後は戻り値の型注釈であり、実行時の値を自動変換する命令ではない。
\item 内包表記は、反復・条件・値生成を一つの式で記述してcollectionを作る。
\item with文はcontext managerに開始・終了処理を任せ、ファイルなどの資源を確実に解放する。
            \end{itemize}
            \paragraph{上から順の処理}
            \begin{enumerate}[leftmargin=1.8em]
            \item resolved に Path(path).expanduser().resolve() の結果を代入する。
\item 条件 not resolved.is\_file() を判定し、真なら内部処理を行う。
\item   FileNotFoundError(f'\{label\} CSV was not found: \{resolved\}') を送出する。
\item with 文で resolved.open('r', encoding='utf-8-sig', newline='') を管理しながら処理する。
\item   reader に csv.DictReader(stream) の結果を代入する。
\item   header に tuple(reader.fieldnames or ()) の結果を代入する。
\item   missing に [column for column in required\_columns if column not in header] の結果を代入する。
\item   条件 missing を判定し、真なら内部処理を行う。
\item     ValueError(f'\{label\} CSV is missing columns \{missing\}: \{resolved\}') を送出する。
\item   条件 next(reader, None) is None を判定し、真なら内部処理を行う。
\item     ValueError(f'\{label\} CSV has a header but no data rows: \{resolved\}. Fill the template or select an identified vehicle profile before launch.') を送出する。
\item resolved を返す。
            \end{enumerate}
            \paragraph{代表コード断片}
            \begin{lstlisting}[language=Python]
def require_csv_data_rows(
    path: str | Path,
    *,
    label: str,
    required_columns: tuple[str, ...] = (),
) -> Path:
    resolved = Path(path).expanduser().resolve()
    if not resolved.is_file():
        raise FileNotFoundError(f"{label} CSV was not found: {resolved}")
    with resolved.open("r", encoding="utf-8-sig", newline="") as stream:
        reader = csv.DictReader(stream)
        header = tuple(reader.fieldnames or ())
        missing = [column for column in required_columns if column not in header]
        if missing:
            raise ValueError(f"{label} CSV is missing columns {missing}: {resolved}")
        if next(reader, None) is None:
            raise ValueError(
                f"{label} CSV has a header but no data rows: {resolved}. "
                "Fill the template or select an identified vehicle profile before launch."
            )
    return resolved
\end{lstlisting}

\subsection{L37 関数 \_resolve\_profile\_path}
            \begin{itemize}[leftmargin=1.6em]
              \item 行範囲: L37--L48
              \item 所属: module直下
              \item 定義: \path!_resolve_profile_path(path_like: str | Path) -> Path!
              \item docstring: 明示されていない。
              \item このブロックが直接呼ぶ主な関数・メソッド: Path, cwd, exists, expanduser, is\_absolute, resolve, str, strip
              \item 戻り値の要点: candidates[-1] / raw.resolve() / candidate
              \item 主なローカル代入名: candidate, candidates, raw
              \item 読み取る主なインスタンス属性: 特になし。
              \item 更新する主なインスタンス属性: 特になし。
              \item 明示的に送出する例外: 明示的raiseはない。
              \item 制御構造の規模: 条件分岐 2、ループ 1、try 0
            \end{itemize}
            \paragraph{この定義を読むためのPython構文}
            \begin{itemize}[leftmargin=1.6em]
            \item `->`以後は戻り値の型注釈であり、実行時の値を自動変換する命令ではない。
            \end{itemize}
            \paragraph{上から順の処理}
            \begin{enumerate}[leftmargin=1.8em]
            \item raw に Path(str(path\_like or '').strip()).expanduser() の結果を代入する。
\item 条件 raw.is\_absolute() を判定し、真なら内部処理を行う。
\item   raw.resolve() を返す。
\item candidates に [(Path.cwd() / raw).resolve(), (ROOT / raw).resolve()] の結果を代入する。
\item candidates を順に走査し、各要素を candidate に入れて処理する。
\item   条件 candidate.exists() を判定し、真なら内部処理を行う。
\item     candidate を返す。
\item candidates[-1] を返す。
            \end{enumerate}
            \paragraph{代表コード断片}
            \begin{lstlisting}[language=Python]
def _resolve_profile_path(path_like: str | Path) -> Path:
    raw = Path(str(path_like or "").strip()).expanduser()
    if raw.is_absolute():
        return raw.resolve()
    candidates = [
        (Path.cwd() / raw).resolve(),
        (ROOT / raw).resolve(),
    ]
    for candidate in candidates:
        if candidate.exists():
            return candidate
    return candidates[-1]
\end{lstlisting}

\subsection{L51 関数 resolve\_relative\_path}
            \begin{itemize}[leftmargin=1.6em]
              \item 行範囲: L51--L65
              \item 所属: module直下
              \item 定義: \path!resolve_relative_path(base_dir: str | Path, path_like: str | Path) -> str!
              \item docstring: 明示されていない。
              \item このブロックが直接呼ぶ主な関数・メソッド: Path, exists, expanduser, is\_absolute, resolve, str, strip
              \item 戻り値の要点: str(candidate) / '' / str(path.resolve()) / str(candidate)
              \item 主なローカル代入名: base, candidate, path, raw, repo\_candidate
              \item 読み取る主なインスタンス属性: 特になし。
              \item 更新する主なインスタンス属性: 特になし。
              \item 明示的に送出する例外: 明示的raiseはない。
              \item 制御構造の規模: 条件分岐 4、ループ 0、try 0
            \end{itemize}
            \paragraph{この定義を読むためのPython構文}
            \begin{itemize}[leftmargin=1.6em]
            \item `->`以後は戻り値の型注釈であり、実行時の値を自動変換する命令ではない。
            \end{itemize}
            \paragraph{上から順の処理}
            \begin{enumerate}[leftmargin=1.8em]
            \item raw に str(path\_like or '').strip() の結果を代入する。
\item 条件 not raw を判定し、真なら内部処理を行う。
\item   '' を返す。
\item path に Path(raw).expanduser() の結果を代入する。
\item 条件 path.is\_absolute() を判定し、真なら内部処理を行う。
\item   str(path.resolve()) を返す。
\item base に Path(base\_dir).resolve() の結果を代入する。
\item candidate に (base / path).resolve() の結果を代入する。
\item 条件 candidate.exists() を判定し、真なら内部処理を行う。
\item   str(candidate) を返す。
\item repo\_candidate に (ROOT / path).resolve() の結果を代入する。
\item 条件 repo\_candidate.exists() を判定し、真なら内部処理を行う。
\item   str(repo\_candidate) を返す。
\item str(candidate) を返す。
            \end{enumerate}
            \paragraph{代表コード断片}
            \begin{lstlisting}[language=Python]
def resolve_relative_path(base_dir: str | Path, path_like: str | Path) -> str:
    raw = str(path_like or "").strip()
    if not raw:
        return ""
    path = Path(raw).expanduser()
    if path.is_absolute():
        return str(path.resolve())
    base = Path(base_dir).resolve()
    candidate = (base / path).resolve()
    if candidate.exists():
        return str(candidate)
    repo_candidate = (ROOT / path).resolve()
    if repo_candidate.exists():
        return str(repo_candidate)
    return str(candidate)
\end{lstlisting}

\subsection{L68 関数 load\_profile}
            \begin{itemize}[leftmargin=1.6em]
              \item 行範囲: L68--L74
              \item 所属: module直下
              \item 定義: \path!load_profile(path_like: str | Path) -> tuple[Path, dict]!
              \item docstring: 明示されていない。
              \item このブロックが直接呼ぶ主な関数・メソッド: ValueError, \_resolve\_profile\_path, isinstance, open, safe\_load
              \item 戻り値の要点: (profile\_path, cfg)
              \item 主なローカル代入名: cfg, f, profile\_path
              \item 読み取る主なインスタンス属性: 特になし。
              \item 更新する主なインスタンス属性: 特になし。
              \item 明示的に送出する例外: ValueError(f'profile must be a mapping: \{profile\_path\}')
              \item 制御構造の規模: 条件分岐 1、ループ 0、try 0
            \end{itemize}
            \paragraph{この定義を読むためのPython構文}
            \begin{itemize}[leftmargin=1.6em]
            \item `->`以後は戻り値の型注釈であり、実行時の値を自動変換する命令ではない。
\item with文はcontext managerに開始・終了処理を任せ、ファイルなどの資源を確実に解放する。
            \end{itemize}
            \paragraph{上から順の処理}
            \begin{enumerate}[leftmargin=1.8em]
            \item profile\_path に \_resolve\_profile\_path(path\_like) の結果を代入する。
\item with 文で profile\_path.open('r', encoding='utf-8') を管理しながら処理する。
\item   cfg に yaml.safe\_load(f) or \{\} の結果を代入する。
\item 条件 not isinstance(cfg, dict) を判定し、真なら内部処理を行う。
\item   ValueError(f'profile must be a mapping: \{profile\_path\}') を送出する。
\item (profile\_path, cfg) を返す。
            \end{enumerate}
            \paragraph{代表コード断片}
            \begin{lstlisting}[language=Python]
def load_profile(path_like: str | Path) -> tuple[Path, dict]:
    profile_path = _resolve_profile_path(path_like)
    with profile_path.open("r", encoding="utf-8") as f:
        cfg = yaml.safe_load(f) or {}
    if not isinstance(cfg, dict):
        raise ValueError(f"profile must be a mapping: {profile_path}")
    return profile_path, cfg
\end{lstlisting}

\subsection{L77 関数 merged\_dict}
            \begin{itemize}[leftmargin=1.6em]
              \item 行範囲: L77--L90
              \item 所属: module直下
              \item 定義: \path!merged_dict(*payloads: dict | None) -> dict!
              \item docstring: 明示されていない。
              \item このブロックが直接呼ぶ主な関数・メソッド: \_merge, deepcopy, get, isinstance, items
              \item 戻り値の要点: out / dst
              \item 主なローカル代入名: key, out, payload, value
              \item 読み取る主なインスタンス属性: 特になし。
              \item 更新する主なインスタンス属性: 特になし。
              \item 明示的に送出する例外: 明示的raiseはない。
              \item 制御構造の規模: 条件分岐 2、ループ 2、try 0
            \end{itemize}
            \paragraph{この定義を読むためのPython構文}
            \begin{itemize}[leftmargin=1.6em]
            \item `->`以後は戻り値の型注釈であり、実行時の値を自動変換する命令ではない。
            \end{itemize}
            \paragraph{上から順の処理}
            \begin{enumerate}[leftmargin=1.8em]
            \item 関数 \_merge を定義する。
\item out に \{\} を代入する。
\item payloads を順に走査し、各要素を payload に入れて処理する。
\item   条件 isinstance(payload, dict) を判定し、真なら内部処理を行う。
\item     out に \_merge(out, payload) の結果を代入する。
\item out を返す。
            \end{enumerate}
            \paragraph{代表コード断片}
            \begin{lstlisting}[language=Python]
def merged_dict(*payloads: dict | None) -> dict:
    def _merge(dst: dict, src: dict) -> dict:
        for key, value in (src or {}).items():
            if isinstance(value, dict) and isinstance(dst.get(key), dict):
                dst[key] = _merge(copy.deepcopy(dst[key]), value)
            else:
                dst[key] = copy.deepcopy(value)
        return dst

    out: dict[str, Any] = {}
    for payload in payloads:
        if isinstance(payload, dict):
            out = _merge(out, payload)
    return out
\end{lstlisting}

\subsection{L78 関数 merged\_dict.\_merge}
            \begin{itemize}[leftmargin=1.6em]
              \item 行範囲: L78--L84
              \item 所属: \path|merged_dict|
              \item 定義: \path|_merge(dst: dict, src: dict) -> dict|
              \item docstring: 明示されていない。
              \item このブロックが直接呼ぶ主な関数・メソッド: \_merge, deepcopy, get, isinstance, items
              \item 戻り値の要点: dst
              \item 主なローカル代入名: key, value
              \item 読み取る主なインスタンス属性: 特になし。
              \item 更新する主なインスタンス属性: 特になし。
              \item 明示的に送出する例外: 明示的raiseはない。
              \item 制御構造の規模: 条件分岐 1、ループ 1、try 0
            \end{itemize}
            \paragraph{この定義を読むためのPython構文}
            \begin{itemize}[leftmargin=1.6em]
            \item `->`以後は戻り値の型注釈であり、実行時の値を自動変換する命令ではない。
            \end{itemize}
            \paragraph{上から順の処理}
            \begin{enumerate}[leftmargin=1.8em]
            \item (src or \{\}).items() を順に走査し、各要素を (key, value) に入れて処理する。
\item   条件 isinstance(value, dict) and isinstance(dst.get(key), dict) を判定し、真なら内部処理を行う。
\item     dst[key] に \_merge(copy.deepcopy(dst[key]), value) の結果を代入する。
\item     上の条件が偽の場合:
\item     dst[key] に copy.deepcopy(value) の結果を代入する。
\item dst を返す。
            \end{enumerate}
            \paragraph{代表コード断片}
            \begin{lstlisting}[language=Python]
    def _merge(dst: dict, src: dict) -> dict:
        for key, value in (src or {}).items():
            if isinstance(value, dict) and isinstance(dst.get(key), dict):
                dst[key] = _merge(copy.deepcopy(dst[key]), value)
            else:
                dst[key] = copy.deepcopy(value)
        return dst
\end{lstlisting}

\subsection{L93 関数 get\_section}
            \begin{itemize}[leftmargin=1.6em]
              \item 行範囲: L93--L105
              \item 所属: module直下
              \item 定義: \path|get_section(cfg: dict, key: str, default: Any = None) -> Any|
              \item docstring: 明示されていない。
              \item このブロックが直接呼ぶ主な関数・メソッド: deepcopy, get, isinstance, split, str
              \item 戻り値の要点: copy.deepcopy(current) / copy.deepcopy(default) if default is not None else \{\} / copy.deepcopy(default) if default is not None else \{\} / copy.deepcopy(default) if default is not None else \{\}
              \item 主なローカル代入名: current, part
              \item 読み取る主なインスタンス属性: 特になし。
              \item 更新する主なインスタンス属性: 特になし。
              \item 明示的に送出する例外: 明示的raiseはない。
              \item 制御構造の規模: 条件分岐 4、ループ 1、try 0
            \end{itemize}
            \paragraph{この定義を読むためのPython構文}
            \begin{itemize}[leftmargin=1.6em]
            \item `->`以後は戻り値の型注釈であり、実行時の値を自動変換する命令ではない。
            \end{itemize}
            \paragraph{上から順の処理}
            \begin{enumerate}[leftmargin=1.8em]
            \item current に cfg を代入する。
\item str(key or '').split('.') を順に走査し、各要素を part に入れて処理する。
\item   条件 not part を判定し、真なら内部処理を行う。
\item     Continue 文を実行する。
\item   条件 not isinstance(current, dict) を判定し、真なら内部処理を行う。
\item     copy.deepcopy(default) if default is not None else \{\} を返す。
\item   current に current.get(part) の結果を代入する。
\item   条件 current is None を判定し、真なら内部処理を行う。
\item     copy.deepcopy(default) if default is not None else \{\} を返す。
\item 条件 current is None を判定し、真なら内部処理を行う。
\item   copy.deepcopy(default) if default is not None else \{\} を返す。
\item copy.deepcopy(current) を返す。
            \end{enumerate}
            \paragraph{代表コード断片}
            \begin{lstlisting}[language=Python]
def get_section(cfg: dict, key: str, default: Any = None) -> Any:
    current: Any = cfg
    for part in str(key or "").split("."):
        if not part:
            continue
        if not isinstance(current, dict):
            return copy.deepcopy(default) if default is not None else {}
        current = current.get(part)
        if current is None:
            return copy.deepcopy(default) if default is not None else {}
    if current is None:
        return copy.deepcopy(default) if default is not None else {}
    return copy.deepcopy(current)
\end{lstlisting}

\subsection{L108 関数 get\_path}
            \begin{itemize}[leftmargin=1.6em]
              \item 行範囲: L108--L118
              \item 所属: module直下
              \item 定義: \path!get_path(cfg: dict, profile_path: str | Path, key: str, default: str = '') -> str!
              \item docstring: 明示されていない。
              \item このブロックが直接呼ぶ主な関数・メソッド: Path, get, get\_section, isinstance, resolve, resolve\_relative\_path, str, strip
              \item 戻り値の要点: resolve\_relative\_path(profile\_dir, raw)
              \item 主なローカル代入名: paths\_cfg, profile\_dir, raw
              \item 読み取る主なインスタンス属性: 特になし。
              \item 更新する主なインスタンス属性: 特になし。
              \item 明示的に送出する例外: 明示的raiseはない。
              \item 制御構造の規模: 条件分岐 3、ループ 0、try 0
            \end{itemize}
            \paragraph{この定義を読むためのPython構文}
            \begin{itemize}[leftmargin=1.6em]
            \item `->`以後は戻り値の型注釈であり、実行時の値を自動変換する命令ではない。
            \end{itemize}
            \paragraph{上から順の処理}
            \begin{enumerate}[leftmargin=1.8em]
            \item profile\_dir に Path(profile\_path).resolve().parent の結果を代入する。
\item raw に '' の結果を代入する。
\item paths\_cfg に cfg.get('paths', \{\}) if isinstance(cfg, dict) else \{\} の結果を代入する。
\item 条件 isinstance(paths\_cfg, dict) を判定し、真なら内部処理を行う。
\item   raw に str(paths\_cfg.get(key, '') or '').strip() の結果を代入する。
\item 条件 not raw を判定し、真なら内部処理を行う。
\item   raw に str(get\_section(cfg, key, default) or '').strip() の結果を代入する。
\item 条件 not raw を判定し、真なら内部処理を行う。
\item   raw に str(default or '').strip() の結果を代入する。
\item resolve\_relative\_path(profile\_dir, raw) を返す。
            \end{enumerate}
            \paragraph{代表コード断片}
            \begin{lstlisting}[language=Python]
def get_path(cfg: dict, profile_path: str | Path, key: str, default: str = "") -> str:
    profile_dir = Path(profile_path).resolve().parent
    raw = ""
    paths_cfg = cfg.get("paths", {}) if isinstance(cfg, dict) else {}
    if isinstance(paths_cfg, dict):
        raw = str(paths_cfg.get(key, "") or "").strip()
    if not raw:
        raw = str(get_section(cfg, key, default) or "").strip()
    if not raw:
        raw = str(default or "").strip()
    return resolve_relative_path(profile_dir, raw)
\end{lstlisting}











        \section{処理の流れ}
        \begin{enumerate}[leftmargin=1.7em]
        \item ソース中の主要関数を通じて処理を進める。
        \end{enumerate}

        \section{このファイルの位置づけ}
        この資料群では、\path|mpc_solarcar/solar_profile.py| を
        \textbf{config}の中核として扱う。
        したがって、ここで扱う処理は単独で完結せず、
        前段の profile / launch / telemetry と後段の planner / logger / report へ繋がる。

        \end{document}
